\begin{proof}
我们给出详细证明，因为它将启发下一小节中 $S$-算子的构造。为简化记号，以下所有公式与图表中都省略结合约束。

我们先定义范畴 $\Tr'_\sigma(\mC)$，其对象与 $\mC$ 的对象相同。为避免混淆，当 $X\in\mC$ 被看作 $\Tr'_\sigma(\mC)$ 的对象时，记作 $\widetilde{X}$。态射空间定义为
\begin{equation*}
\label{E:trace hom}
\Hom_{\Tr'_\sigma(\mC)}(\widetilde{X},\widetilde{Y}):=\big(\bigoplus_{(V,W)\in \mC\times \mC}\Hom(X, V\otimes W)\otimes_E \Hom(W\otimes \sigma V, Y)\big)/\sim,
\end{equation*}
其中等价关系由下述方式生成：若 $u: X\to V\otimes W$ 与 $v:W\otimes \sigma V\to Y$ 对应的元素记作 $S_{u,v}\in \Hom_{\Tr'_\sigma(\mC)}(\widetilde{X},\widetilde{Y})$，则当存在交换图
\[\xymatrix{
X\ar^-{u}[r] \ar_{u'}[dr]& V\otimes W\ar[d] & W\otimes \sigma V\ar[d]\ar^-{v}[r] & Y\\
& V'\otimes W' & W'\otimes \sigma V'\ar_{v'}[ur] &
}\]
（其中竖直映射由 $\Hom_\mC(V,V')\otimes_E\Hom_{\mC}(W,W')$ 中的元素诱导）时，令 $S_{u,v}=S_{u',v'}$。特别地，由下面两个交换图
\begin{equation}
\label{E:frac}
\xymatrix{
&V\otimes V^*\otimes X\ar^{\id_V\otimes u^*}[d]& V^*\otimes X\otimes \sigma V\ar^{v_u}[dr]\ar_{u^*\otimes \id_{\sigma V}}[d]&\\
X\ar^{\coev_V\otimes\id_X}[ur]\ar^{u}[r]\ar_{u_v}[dr]&V\otimes W\ar^{\id_V\otimes v^*}[d]&W\otimes \sigma V\ar^{v}[r]\ar_{v^*\otimes\id_{\sigma V}}[d]& Y\\
&V\otimes Y\otimes \sigma V^*& Y\otimes \sigma V^*\otimes \sigma V\ar_{\id_Y\otimes\ev_{\sigma V}}[ur]&
}
\end{equation}
可得 $S_{\coev_V\otimes\id_X, v_u}=S_{u,v}=S_{u_v,\id_Y\otimes\ev_{\sigma V}}$。
这里 $V^*$ 是 $V$ 的左对偶，$u^*: V^*\otimes X\to W$ 是 $u$ 的对偶，$v^*:W\to Y\otimes \sigma V^*$ 是 $v$ 的对偶，而
$u_v= (\id_V\otimes v^*)\circ u$，$v_u= v\circ(u^*\otimes \id_{\sigma V})$。

我们还需要说明 $\Tr'_\sigma(\mC)$ 中态射的复合。给定 $u_1: X\to V_1\otimes W_1$ 与 $v_1: W_1\otimes \sigma V_1\to Y$，再给定 $u_2: Y\to V_2\otimes W_2$ 与 $v_2: W_2\otimes \sigma V_2\to Z$，设 $S_{u_1,v_1}\in \Hom_{\Tr'_\sigma(\mC)}(\widetilde{X},\widetilde{Y})$ 与 $S_{u_2,v_2}\in \Hom_{\Tr'_\sigma(\mC)}(\widetilde{Y},\widetilde{Z})$ 是对应的元素。我们定义 $S_{u_2,v_2}\circ S_{u_1,v_1}$ 为 $S_{u,v}$，其中
\[u: X\xrightarrow{u_1} V_1\otimes W_1\xrightarrow{\id_{V_1}\otimes\coev_{V_2}\otimes\id_{W_1}}V_1\otimes V_2\otimes V^*_2\otimes W_1,\]
\[v: V^*_2\otimes W_1\otimes \sigma V_1\otimes \sigma V_2\xrightarrow{\id_{V^*_2}\otimes v_1\otimes\id_{\sigma V_2}}V^*_2\otimes Y\otimes \sigma V_2\xrightarrow{(v_2)_{u_2}} Z.\]

立即可验证该复合与 $(u_1,v_1)$ 的选取无关。另一方面，利用 \eqref{E:frac} 容易看出 $S_{u,v}=S_{u',v'}$，其中
\[
u': X\xrightarrow{(u_1)_{v_1}} V_1\otimes Y\otimes \sigma V^*_1\xrightarrow{\id_{V_1}\otimes u_2\otimes \id_{\sigma V^*_1}} V_1\otimes V_2\otimes W_2\otimes \sigma V^*_1,
\]
\[
v': W_2\otimes \sigma V^*_1 \otimes \sigma V_1\otimes \sigma V_2 \xrightarrow{\id_{W_2}\otimes\ev_{\sigma V_1}\otimes\id_{\sigma V_2}}W_2\otimes \sigma V_2\xrightarrow{v_2} Z.
\]

由这一表达式可验证该复合同样与 $(u_2,v_2)$ 的选取无关。
此外，立即可验证 $\widetilde{X}$ 的恒等态射可由典范同构 $X\cong  \mathbf{1}\otimes X$ 与 $X\otimes \sigma \mathbf{1}\cong X\otimes\mathbf{1}\cong X$ 代表。

至此我们已把 $\Tr'_\sigma(\mC)$ 定义成一个范畴。给定 $X$ 与 $Y$，有典范态射
\[\al_{X,Y}:\widetilde{X\otimes Y}\to \widetilde{Y\otimes \sigma X},\]
它由 $S_{u,v}$ 给出，其中 $u$ 与 $v$ 都取恒等映射。
特别地，对每个 $X$ 都有典范态射
$$\ga_{X}:\widetilde{X}\cong \widetilde{X\otimes\mathbf{1}}\xrightarrow{\al_{X,\mathbf{1}}} \widetilde{\mathbf{1}\otimes \sigma X}\cong \widetilde{\sigma X}.$$
由构造可验证族 $\{\al_{X,Y}\}$ 满足交换图 \eqref{E:cond trace}，从而 \eqref{E:cond trace cor} 也成立。
我们把 $\Tr_\sigma(\mC)$ 定义为 $\Tr'_\sigma(\mC)$ 关于乘法系 $\{\ga^n_X,\ X\in \mC, n\in \bZ_{\geq 0}\}$ 的局部化（于是 $\ga_X$ 在 $\Tr_\sigma(\mC)$ 中成为同构）。由此所有 $\al_{X,Y}$ 在 $\Tr_\sigma(\mC)$ 中都成为同构。
\end{proof}

\begin{rmk}
定义复合 $S_{u_2,v_2}\circ S_{u_1,v_1}$ 的想法可由下图说明
\begin{tiny}
\begin{equation*}
\label{E: comp in tr}
\xymatrix@C=0.8em@R=0.9em{
                     &                                                                                          &V_1\otimes V_2\otimes V^*_2\otimes W_1\ar@{-->}[r]&V_2\otimes V_2^*\otimes W_1\otimes \sigma V_1\ar[dr]\ar@{-->}[r]& V_2^*\otimes W_1\otimes \sigma V_1\otimes \sigma V_2\ar^-{\id_{V^*_2}\otimes v_1\otimes\id_{\sigma V_2}}[dr] &                                                    &\\
 X\ar^-{u_1}[r]&V_1\otimes W_1\ar^-{\id\otimes\coev_{V_2}\otimes\id}[ur]\ar@{-->}[r]& W_1\otimes \sigma V_1\ar[ur]\ar[r]                      &                 Y\ar[r]                                                   & V_2\otimes V_2^*\otimes Y    \ar@{-->}[r]                              & V_2^*\otimes Y\otimes \sigma V_2 \ar^-{(v_2)_{u_2}}[r]& Z. \\
}
\end{equation*}
\end{tiny}

也就是说，我们可以先把 $S_{u_2,v_2}$ 换成 $S_{\coev_{V_2}\otimes \id_Y,(v_2)_{u_2}}$。由于中间的三角形在 $\mC$ 中交换，且虚线箭头在 $\Tr_\sigma(\mC)$ 中应当是同构，$S_{\coev_{V_2}\otimes \id_Y,(v_2)_{u_2}}\circ S_{u_1,v_1}$ 必定就是证明中给出的那个元素。
\end{rmk}

\begin{rmk}
当然，还有右迹函子的对偶概念，即一个函子 $F:\mC\to \mD$ 连同函子性同构族 $\beta_{X,Y}: F(X\otimes Y)\cong F(\sigma Y\otimes X)$，满足与 \eqref{E:cond trace} 类似的条件；相应地有右 $\sigma$-扭迹范畴的概念。上述构造与下文的讨论都有对偶的版本。

为简化叙述，以下我们把左迹函子简称迹函子，把 $\mC$ 的左 $\sigma$-扭迹范畴简称 $\mC$ 的 $\sigma$-扭迹范畴。
\end{rmk}

下面是上述构造的一个具体例子。

\begin{ex}
\label{E: trace of rep}
设 $H$ 是特征零域 $E$ 上的分裂约化代数群，带有自同构 $\sigma: H\to H$。考虑由公式 \eqref{E:twist conj} 给出的 $H$ 在自身上的 $\sigma$-扭共轭，它等价于 $H$ 在旁集 $H\sigma\subset H\rtimes\langle\sigma\rangle$ 上的通常共轭。记 $\frac{H}{c_\sigma H}$ 为相应的商叠。
我们把 $H$ 的有限维表示范畴 $\on{Rep}_E(H)$ 看作分类叠 $\mathbf BH$ 上凝聚层的范畴。则拉回函子 $\pi^*: \on{Rep}_E(H)\to \Coh(\frac{H}{c_\sigma H})$ 带有 $\sigma$-扭迹函子结构，它诱导
\[\Tr_\sigma(\on{Rep}_E(H))\cong \on{Coh}_{fr}^H(H\sigma)\subset \on{Coh}(\frac{H}{c_\sigma H}),\]
其中 $\Coh_{fr}^H(H\sigma)$ 表示 $\on{Coh}(\frac{H}{c_\sigma H})$ 中由这些 $\widetilde V(:=\pi^*V)$ 张成的全子范畴。这可以这样看出。
经由下降，可以把 $\widetilde V$ 看作 $H$ 上的平凡向量丛 $H\times V$，配以对角的 $H$-等变结构。此外，对表示 $H\to \GL(V)$，记 $\sigma V$ 为 $V$ 的 $\sigma$-扭，即由 $H\xrightarrow{\sigma^{-1}}H\to\GL(V)$ 定义的 $H$ 的表示。我们形式地把恒等映射 $V\to \sigma V$ 记作 $\sigma$，于是 $h\cdot \sigma=\sigma\cdot \sigma^{-1}(h):V\to\sigma V$。
则对 $\Tr'_\sigma(\on{Rep}_E(H))$ 中的态射 $S_{u,v}$（由 $u:X\to V\otimes W$ 与 $v:W\otimes \sigma V\to Y$ 给出），我们定义 $\Coh_{fr}^H(H\sigma)$ 中的态射 $\widetilde{X}\to\widetilde{Y}$，它在 $h\in H$ 上纤维的限制为
\[
\widetilde{X}|_h=X\xrightarrow{u} V\otimes W\xrightarrow{h\sigma\otimes\id}\sigma V\otimes W\cong W\otimes \sigma V\xrightarrow{v} Y=\widetilde{Y}|_h.
\]
这给出自然函子 $\Tr_\sigma(\on{Rep}_E(H))\to \Coh_{fr}^H(H\sigma)$。按定义它是本质满的，而 Peter–Weyl 定理说明它也是全忠实的，因此是同构。

设 $\sigma=\id$。注意由于 $\on{Rep}_E(H)$ 是对称幺半的，恒等函子是迹函子，因此诱导函子 $\Tr_\sigma(\on{Rep}_E(H))\cong \Coh_{fr}^H(H\sigma)\to \on{Rep}_E(H)$。可验证它正是 $\frac{H}{cH}$ 上层的限制到 $H$ 的单位元处。
\end{ex}
